\chapter{The informed investor}\label{ch:trading}

\begin{chapterplan}
The question of this chapter: what does an informed order earn when the price already counts the entry that the
order brings, and which features of the auction does the answer use?

After this chapter you can:
\begin{itemize}
\item derive the gross return of a purchase in class \(H\) and of a sale in class \(L\), one line at a time;
\item explain why a wrong-signed order never pays, and when each class buys, sells or abstains;
\item derive the trading bound and explain why the class that the prior favors binds;
\item explain why orders, posteriors, entry and research depend on the auction only through the spread \(\DT\)
and the threshold \(\tau\);
\item show that a stronger incumbent never narrows the spread and, while \(g_L\le c_H\), never lowers the
threshold.
\end{itemize}
\end{chapterplan}

% =====================================================================
\section{What an informed order can earn}\label{sec:trading-intuition}

The informed investor knows the challenger's class. When it fills the order slot, it can buy one unit of the
target's stock, sell one unit, or place no order. The market makers see the order but not who placed it. They set
the price equal to the expected terminal payoff of a share, given the order. The challenger then reads the price
and decides whether to prepare.

Take a purchase by the investor in class \(H\). The purchase moves the market makers' belief toward class \(H\).
If the belief rises far enough, a high-cost challenger prepares, so entry rises and expected target proceeds rise
with it. The market makers know how the challenger responds to each price, so they put the expected entry gain
into the price. The investor therefore pays, inside the price, for the entry that its own order brings.

What is left for the investor is the one thing that the market makers do not know: the class. After a purchase,
the market makers still give some weight to class \(L\), because a liquidity trader can place the same
order. They value the share as an average over the two classes, and the investor values it at the class-\(H\)
level. The gap is the spread \(\DT\) times the weight that the market gives to the wrong class. This gap matters
only when the challenger prepares, because without entry the target receives \(t_0\) in both classes.

\begin{keyidea}
The gross return of an informed order is entry times the spread times the residual uncertainty about the class.
For a purchase in class \(H\) it is \(e(1)\,\DT\,[1-\mu(1)]\). For a sale in class \(L\) it is
\(e(-1)\,\DT\,\mu(-1)\).
\end{keyidea}

\begin{analogy}[a bet on a match that may not take place]
This analogy is the guide's own. A ticket on a football match pays a fixed base amount if the match does not take
place. If the match takes place, the ticket pays more on average when a star player plays. You know whether the
star plays. A bookmaker sees your order, but not whether you know anything, and sets a fair price. The fair price
contains the base amount, the chance of the match, and the bookmaker's belief about the star. So the base amount
and the chance of the match give you no edge. Your edge is the chance of the match, times the payout gap, times
the bookmaker's doubt about the star.

\emph{Where the analogy breaks.} In the model, the price itself changes the chance that the ``match'' takes place,
because the challenger reads the price before it prepares. A bookmaker's price does not decide whether a match
takes place. The analogy fits the gross return at fixed prices and entry, which is how the investor evaluates an
order. It says nothing about how entry responds to prices in equilibrium.
\end{analogy}

Each of the three factors can remove the return. If the class did not affect proceeds (\(\DT=0\)), knowledge of
the class would be worth nothing. If a purchase revealed the class completely, the price would already contain
everything that the investor knows. If the challenger never prepared, a share would pay \(t_0\) in both classes.
The rest of the chapter makes this story exact.

% =====================================================================
\section{The set-up}\label{sec:trading-setup}

This section recalls the symbols that the chapter uses.
\begin{itemize}
\item The challenger's class is \(j\in\{H,L\}\), with prior \(\pi=\Pr(j=H)\in(0,1)\).
\item With probability \(\lambda\in(0,1)\) the informed investor fills the order slot. It sees the class and
chooses an order \(q\in\{-1,0,1\}\): sell, abstain, or buy one unit. A nonzero order costs \(k>0\).
\item With probability \(1-\lambda\) a liquidity trader fills the slot and sells, abstains or buys with
probabilities \(\beta\), \(1-2\beta\) and \(\beta\). A liquidity order in one direction arrives with probability
\(b=(1-\lambda)\beta\), and \(u=\lambda/b\).
\item After the order \(x\), the market makers' posterior is \(\mu(x)=\Pr(j=H\mid X=x)\), and the price is \(P(x)\).
\item Entry \(e(x)\) is the probability that the challenger prepares after the price \(P(x)\). A zero-cost
challenger appears with probability \(\rho\in(0,1)\) and prepares after every price, so
\(e(x)=\rho+(1-\rho)\,\zeta(P(x))\in[\rho,1]\), where \(\zeta\) is the probability that a high-cost challenger
prepares.
\item Expected target proceeds are \(t_0\) without entry and \(t_j\) after entry in class \(j\). The target-payoff
spread is \(\DT=t_H-t_L\).
\end{itemize}

The paper's appendix uses two short names. The \emph{entry gain} in class \(j\) is \(w_j=t_j-t_0\), the rise in
expected proceeds that entry brings. Its average at the belief \(\mu\) is \(w(\mu)=(1-\mu)\,w_L+\mu\,w_H\).
Proposition~1 and Chapter~\ref{ch:auction} give \(w_H\ge w_L\ge0\). Section~\ref{sec:beliefs-setup} derives the short form of the average, \(w(\mu)=w_L+\DT\,\mu=t_L-t_0+\DT\,\mu\), because \(w_H-w_L=t_H-t_L=\DT\).

\paragraph{The investor's objective.} If the investor places the order \(q\), its order is the whole order that the
market makers see. It trades at the price \(P(q)\) and earns \(q\,[V_T-P(q)]-k|q|\), where \(V_T\) is the terminal
payoff of a share. Section~\ref{sec:model-equilibrium} of Chapter~\ref{ch:model} takes the expectation given the class \(j\) and
the order \(q\):
\begin{equation}\label{eq:trading-objective}
U_j(q)=q\,\bigl[t_0+e(q)\,(t_j-t_0)-P(q)\bigr]-k|q| .
\end{equation}
The first term in the bracket is the expected terminal payoff given the class and the order. Steps 1 to 3 of
Section~\ref{sec:beliefs-priceformula} derive it, in \eqref{eq:beliefs-terminal}:
\[
\E[V_T\mid j,X=q]=t_0+e(q)\,w_j .
\]
In words: the challenger prepares with probability \(e(q)\) in both classes, because its choice depends only on the
price and its cost. If it prepares, the target expects \(t_j\); if not, it expects \(t_0\).

\paragraph{The price.} The market makers set \(P(x)=\E[V_T\mid X=x]\). Step 4 of
Section~\ref{sec:beliefs-priceformula} averages \eqref{eq:beliefs-terminal} over the class, with the weights
\(\mu(x)\) and \(1-\mu(x)\), and obtains the price formula \eqref{eq:beliefs-price}:
\[
P(x)=\mu(x)\,[t_0+e(x)\,w_H]+[1-\mu(x)]\,[t_0+e(x)\,w_L]=t_0+e(x)\,\bigl[t_L-t_0+\DT\,\mu(x)\bigr].
\]
The first form shows the fact behind every cancellation in this chapter. The price is the market makers' average,
over the two classes, of the same expected payoff that enters the investor's objective. The investor in class \(H\)
receives the class-\(H\) version and pays the average.

\begin{watchout}
The price formula needs only Bayes consistency (E1), competitive prices (E4), truthful bids (E6) and a cost that is
independent of the class. It does not need a positive spread. What needs a positive spread \emph{and} a positive
entry floor \(\rho\) is the step of Section~\ref{sec:beliefs-reveal}: that the price reveals the posterior. Also,
the price is affine in the posterior only at a fixed level of entry. Equilibrium prices reflect entry too.
\end{watchout}

\paragraph{Deviations and the posterior bounds.} When the investor considers the order \(q\), it takes the
candidate's price \(P(q)\) and entry \(e(q)\) as given. A comparison across economies, such as across incumbent
strengths, is a different exercise: it solves the equilibrium again in each economy. Recall the posterior bounds
\eqref{eq:beliefs-bounds}:
\[
M(\pi)=\frac{\pi(1+u)}{\pi(1+u)+1-\pi},\qquad m(\pi)=\frac{\pi}{\pi+(1-\pi)(1+u)} .
\]
Under every investor strategy, mixed or pure, the posterior after a purchase or a sale lies between these two
bounds, for every \(0<\beta<1/2\) (Lemma A.1(b)). The reason is that each class reaches each nonzero order through
a liquidity trader with probability \(b>0\).

% =====================================================================
\section{The gross return of an order}\label{sec:trading-returns}

The paper collects the investor's problem in one lemma. This section proves its six payoffs and part (a), and
Sections~\ref{sec:trading-best} to~\ref{sec:trading-split} prove parts (b) to (d).

\begin{paperresult}{Lemma A.3 (Informed trading)}
Let \(\DT\ge0\). Fix a tuple that satisfies (E1), (E4) and (E6), with entry \(e(x)\in[\rho,1]\) after every order.
Conditional on its arrival, the investor in each class earns
\[
\begin{aligned}
U_H(1)&=e(1)\DT[1-\mu(1)]-k,& U_H(0)&=0,& U_H(-1)&=-e(-1)\DT[1-\mu(-1)]-k,\\
U_L(-1)&=e(-1)\DT\,\mu(-1)-k,& U_L(0)&=0,& U_L(1)&=-e(1)\DT\,\mu(1)-k.
\end{aligned}
\]
Write \(K_H=e(1)\DT[1-\mu(1)]\) and \(K_L=e(-1)\DT\,\mu(-1)\) for the gross returns of the two correct orders.
\begin{enumerate}[label=(\alph*)]
\item A wrong-signed order earns at most \(-k\), so abstention strictly beats it.
\item Class \(H\) buys if \(K_H>k\) and abstains if \(K_H<k\). At \(K_H=k\), the purchase and abstention are its two
best orders. Class \(L\) sells if \(K_L>k\) and abstains if \(K_L<k\). At \(K_L=k\), the sale and abstention are
its two best orders.
\item Under every candidate profile, \(\rho\DT[1-M(\pi)]\le K_H\le\DT\) and \(\rho\DT\,m(\pi)\le K_L\le\DT\).
Hence \(\rho\,\kappa(\pi)\DT\le K_j\le\DT\), where \(\kappa(\pi)=\min\{m(\pi),1-M(\pi)\}\).
\item \(K_H\) depends on the candidate only through \(\mu(1)\) and \(e(1)\), and \(K_L\) only through \(\mu(-1)\)
and \(e(-1)\).
\end{enumerate}
\end{paperresult}

A plain reading:
\begin{itemize}
\item \emph{Hypotheses.} The spread can be zero. The tuple needs Bayes consistency, competitive prices and truthful
bids, and entry between \(\rho\) and one. It need not be an equilibrium, so the lemma holds in and out of
equilibrium.
\item \emph{Payoffs.} ``Conditional on its arrival'' means: given that the investor fills the slot. Each class has
a correct order (a purchase for \(H\), a sale for \(L\)), a wrong-signed order, and abstention, which earns zero.
\(K_H\) and \(K_L\) are the payoffs of the correct orders before the cost.
\item \emph{(a)} The investor never trades against its own information. \emph{(b)} Each class compares its gross
return with the cost. \emph{(c)} The gross returns have a floor, a fixed share of the spread, and a ceiling, the
spread, in every \emph{candidate}: a tuple that meets all equilibrium conditions except, possibly, the investor's
optimality (E5). \emph{(d)} The return of a purchase uses only what happens after a purchase, and the return of a
sale only what happens after a sale.
\end{itemize}

\subsection{Step 1: the gross returns}

Substitute the price formula into the bracket of \eqref{eq:trading-objective}. Start with class \(H\).

\begin{derivation}[the bracket in class \(H\)]
\begin{align*}
&t_0+e(q)(t_H-t_0)-P(q)\\
&\quad= t_0+e(q)(t_H-t_0)-\bigl\{t_0+e(q)\,[t_L-t_0+\DT\,\mu(q)]\bigr\} \why{substitute the price formula}\\
&\quad= e(q)(t_H-t_0)-e(q)\,[t_L-t_0+\DT\,\mu(q)] \why{\(t_0-t_0=0\)}\\
&\quad= e(q)\,\bigl[(t_H-t_0)-(t_L-t_0)-\DT\,\mu(q)\bigr] \why{factor out \(e(q)\)}\\
&\quad= e(q)\,\bigl[t_H-t_L-\DT\,\mu(q)\bigr] \why{\(-t_0+t_0=0\)}\\
&\quad= e(q)\,\bigl[\DT-\DT\,\mu(q)\bigr] \why{\(\DT=t_H-t_L\)}\\
&\quad= e(q)\,\DT\,[1-\mu(q)] \why{factor out \(\DT\)}
\end{align*}
\end{derivation}

\begin{derivation}[the bracket in class \(L\): the same steps with \(t_L\) in place of \(t_H\)]
\begin{align*}
&t_0+e(q)(t_L-t_0)-P(q)\\
&\quad= e(q)(t_L-t_0)-e(q)\,[t_L-t_0+\DT\,\mu(q)] \why{substitute; \(t_0-t_0=0\)}\\
&\quad= e(q)\,\bigl[(t_L-t_0)-(t_L-t_0)-\DT\,\mu(q)\bigr] \why{factor out \(e(q)\)}\\
&\quad= -e(q)\,\DT\,\mu(q) \why{\((t_L-t_0)-(t_L-t_0)=0\)}
\end{align*}
\end{derivation}

Multiply each bracket by \(q\) and subtract \(k|q|\). In class \(H\), \(U_H(q)=q\,e(q)\,\DT\,[1-\mu(q)]-k|q|\); in
class \(L\), \(U_L(q)=-q\,e(q)\,\DT\,\mu(q)-k|q|\).

\begin{derivation}[the six payoffs]
\begin{align*}
U_H(1) &= e(1)\,\DT\,[1-\mu(1)]-k \why{\(q=1\), \(|q|=1\)}\\
U_H(0) &= 0\cdot e(0)\,\DT\,[1-\mu(0)]-0=0 \why{\(q=0\), \(|q|=0\)}\\
U_H(-1) &= -e(-1)\,\DT\,[1-\mu(-1)]-k \why{\(q=-1\), \(|q|=1\)}\\
U_L(-1) &= e(-1)\,\DT\,\mu(-1)-k \why{\(q=-1\): \(-q=1\), \(|q|=1\)}\\
U_L(0) &= 0 \why{\(q=0\)}\\
U_L(1) &= -e(1)\,\DT\,\mu(1)-k \why{\(q=1\), \(|q|=1\)}
\end{align*}
\end{derivation}

These are the six payoffs of Lemma A.3. The two correct orders have the gross returns
\begin{equation}\label{eq:trading-returns}
K_H=e(1)\,\DT\,\bigl[1-\mu(1)\bigr],\qquad K_L=e(-1)\,\DT\,\mu(-1).
\end{equation}

\begin{intuition}[why \(t_0\) cancels]
The proceeds \(t_0\) are what the target receives when no challenger prepares. They enter the terminal payoff and,
with the same coefficient one, the price. The investor receives \(t_0\) at the end and pays \(t_0\) at the start,
so they net out. The low-class entry gain \(e(q)\,w_L\) cancels for the same reason: every class brings at least
\(w_L\) when the challenger prepares, and the market knows it. What remains is the part of the entry gain that
depends on the class, \(\DT\), times the weight that the price gives to the wrong class.
\end{intuition}

The same result reads well in entry gains. The investor in class \(H\) receives its own gain \(w_H\) per unit of
entry and pays the market's average gain \(w(\mu)\):
\begin{align*}
\E[V_T\mid H,X=q]-P(q) &= e(q)\,\bigl[w_H-(1-\mu(q))\,w_L-\mu(q)\,w_H\bigr] \why{\(t_0\) cancels; definition of \(w\)}\\
&= e(q)\,[1-\mu(q)]\,(w_H-w_L) \why{collect the terms}\\
&= e(q)\,\DT\,[1-\mu(q)] \why{\(w_H-w_L=\DT\)}
\end{align*}
In the same way, \(P(q)-\E[V_T\mid L,X=q]=e(q)\,\mu(q)\,(w_H-w_L)=e(q)\,\DT\,\mu(q)\).

\paragraph{What this says.} A purchase by class \(H\) earns, before its cost, entry times the spread times the
market's residual weight on class \(L\). A sale by class \(L\) earns entry times the spread times the residual
weight on class \(H\). Test the extremes.
\begin{itemize}
\item If \(\DT=0\), then \(K_H=K_L=0\), and every nonzero order loses its cost \(k\).
\item A purchase that revealed the class, \(\mu(1)=1\), would give \(K_H=0\). The slot prevents this: a liquidity purchase is
always possible, so \(\mu(1)\le M(\pi)<1\).
\item A purchase that does not recruit has \(e(1)=\rho\). At a fixed posterior, a recruitment multiplies the return
by \(1/\rho\).
\item When the investor is rare (\(u\) small), posteriors stay near the prior, and \(K_H\) is near
\(e(1)\,\DT\,(1-\pi)\).
\end{itemize}

\begin{checkyourself}\label{q:trading-level}
Suppose a change in the auction raises \(t_0\), \(t_H\) and \(t_L\) by the same amount \(z>0\), and leaves
posteriors and entry unchanged. Show, step by step, what happens to each price and to \(K_H\) and \(K_L\).
\end{checkyourself}

\subsection{Step 2: wrong-signed orders}

A wrong-signed order is a sale by class \(H\) or a purchase by class \(L\). Part (a) says that it never pays.

\begin{derivation}[each wrong-signed order earns at most \(-k\)]
For a sale by class \(H\):
\begin{align*}
e(-1)\,\DT\,[1-\mu(-1)] &\ge 0 \why{\(e\ge\rho>0\), \(\DT\ge0\), \(\mu(-1)\le1\)}\\
-e(-1)\,\DT\,[1-\mu(-1)] &\le 0 \why{multiply by \(-1\); the inequality turns}\\
U_H(-1)=-e(-1)\,\DT\,[1-\mu(-1)]-k &\le -k<0=U_H(0) \why{subtract \(k\); \(k>0\)}
\end{align*}
For a purchase by class \(L\):
\begin{align*}
e(1)\,\DT\,\mu(1) &\ge 0 \why{\(e\ge\rho>0\), \(\DT\ge0\), \(\mu(1)\ge0\)}\\
U_L(1)=-e(1)\,\DT\,\mu(1)-k &\le -k<0=U_L(0) \why{multiply by \(-1\) (the inequality turns); subtract \(k\)}
\end{align*}
\end{derivation}

So abstention strictly beats each wrong-signed order, and under the investor's optimality (E5) no class puts
positive probability on one. The definition of equilibrium allows wrong-signed orders; part (a) is what removes
them, for every tuple that meets the hypotheses.

\begin{intuition}[a trade against your own news]
A sale by class \(H\) sells the share at a price that averages in class \(L\). The share is worth more than that to
an investor who knows the class is \(H\). So the sale hands the residual uncertainty to the market makers.
The gross part can be zero, for example when \(\DT=0\), but it is never positive, and the cost \(k\) makes the payoff
strictly negative.
\end{intuition}

% =====================================================================
\section{The best order in each class}\label{sec:trading-best}

By part (a), each class compares only its correct order with abstention. The payoffs of Step 1 give
\[
U_H(1)-U_H(0)=K_H-k,\qquad U_L(-1)-U_L(0)=K_L-k .
\]
This gives part (b) of Lemma A.3, case by case.
\begin{itemize}
\item If \(K_H>k\), the purchase is the unique best order: class \(H\) buys with probability one.
\item If \(K_H<k\), abstention is the unique best order: class \(H\) abstains.
\item If \(K_H=k\), the purchase and abstention are the two best orders, and the investor's tie rule decides. Under
the default tie rules, an indifferent investor may mix over them, that is, buy with any probability. Under the
other investor rule, it abstains.
\end{itemize}
Class \(L\) is the same, with \(K_L\) in place of \(K_H\) and a sale in place of a purchase. The investor's ex ante
payoff is \(\lambda\,U_j(q)\). The factor \(\lambda>0\) multiplies every order's payoff, so it changes no
comparison.

\begin{intuition}[why a mixed order can be a best response]
The return \(K_H\) is not a fixed number. The more often class \(H\) buys, the more a purchase reveals, and the
less residual uncertainty is left to earn on. So at fixed entry the return of a purchase falls as the purchase rate
rises, and a purchase rate can adjust until the return equals the cost. Chapter~\ref{ch:coexistence} uses this.
\end{intuition}

To see the fall, recall that in equilibrium class \(L\) never buys. Write \(s_+=\lambda\,\sigma_H(1)\) for the
probability that the investor in class \(H\) fills the slot and buys. A purchase then arrives with probability
\(b+s_+\) in class \(H\) and \(b\) in class \(L\), and Bayes' rule (Section~\ref{sec:toolkit-bayes}) gives
\begin{align*}
\mu(1) &= \frac{\pi\,(b+s_+)}{\pi\,(b+s_+)+(1-\pi)\,b}=\frac{\pi\,(b+s_+)}{b+\pi s_+} \why{Bayes' rule; \(\pi b+(1-\pi)b=b\)}\\
1-\mu(1) &= \frac{b+\pi s_+-\pi b-\pi s_+}{b+\pi s_+}=\frac{(1-\pi)\,b}{b+\pi s_+} \why{common denominator}
\end{align*}
The numerator does not depend on \(s_+\), and the denominator rises with \(s_+\), so \(1-\mu(1)\) falls. At
\(s_+=0\) it is \(1-\pi\). At \(s_+=\lambda\) it is \((1-\pi)/(1+\pi u)=1-M(\pi)\).

\begin{checkyourself}\label{q:trading-reveal}
Suppose, against the model, that a purchase revealed the class completely, so that \(\mu(1)=1\). What would
\(K_H\) and \(U_H(1)\) be? Which class would then buy? Explain why the order slot rules this situation out.
\end{checkyourself}

% =====================================================================
\section{The trading bound}\label{sec:trading-bound}

Part (c) of Lemma A.3 gives a floor and a ceiling for each gross return. The paper calls it the trading bound:
\begin{equation}\label{eq:trading-bound}
\rho\,\kappa(\pi)\,\DT\;\le\;K_H,\,K_L\;\le\;\DT,\qquad \kappa(\pi)=\min\{m(\pi),\,1-M(\pi)\}.
\end{equation}
Its ingredients are two ranges from earlier chapters: entry lies in \([\rho,1]\), and the posterior after a nonzero
order lies in \([m(\pi),M(\pi)]\).

\begin{derivation}[the floors]
All factors are nonnegative: \(e\ge\rho>0\), \(\DT\ge0\), \(1-\mu(1)\ge1-M(\pi)>0\) and \(\mu(-1)\ge m(\pi)>0\).
So each inequality keeps its direction when it is multiplied by a nonnegative number.
\begin{align*}
K_H &= e(1)\,\DT\,[1-\mu(1)] \why{definition}\\
&\ge \rho\,\DT\,[1-\mu(1)] \why{\(e(1)\ge\rho\); times \(\DT[1-\mu(1)]\ge0\)}\\
&\ge \rho\,\DT\,[1-M(\pi)] \why{\(\mu(1)\le M(\pi)\); times \(\rho\DT\ge0\)}\\
&\ge \rho\,\kappa(\pi)\,\DT \why{\(1-M(\pi)\ge\kappa(\pi)\)}\\[4pt]
K_L &= e(-1)\,\DT\,\mu(-1) \why{definition}\\
&\ge \rho\,\DT\,\mu(-1) \why{\(e(-1)\ge\rho\); times \(\DT\mu(-1)\ge0\)}\\
&\ge \rho\,\DT\,m(\pi) \why{\(\mu(-1)\ge m(\pi)\); times \(\rho\DT\ge0\)}\\
&\ge \rho\,\kappa(\pi)\,\DT \why{\(m(\pi)\ge\kappa(\pi)\)}
\end{align*}
\end{derivation}

For the ceiling, entry and the posterior lie in \([0,1]\), so \(e(1)[1-\mu(1)]\le1\) and \(e(-1)\mu(-1)\le1\).
Multiply by \(\DT\ge0\): \(K_H\le\DT\) and \(K_L\le\DT\). The ceiling uses only \(e\le1\) and \(0\le\mu\le1\).

\paragraph{What this says.} The ceiling would need full entry and a price that puts all its weight on the wrong
class. The posterior bounds rule out the second, so with \(\DT>0\) no order reaches the ceiling.
The floor is the worst case: entry only at the floor \(\rho\), and the posterior pushed as far toward the investor's own
class as the slot allows. Even then the zero-cost challenger prepares, and a liquidity trader hides the order.

\subsection{The shape of \texorpdfstring{\(\kappa(\pi)\)}{kappa}}

The paper writes the bound in a second form, \(\kappa(\pi)=m(\min\{\pi,1-\pi\})\). Two facts give it.

\emph{Fact 1: \(1-M(\pi)=m(1-\pi)\).} This is the mirror identity \eqref{eq:toolkit-mirror} of
Section~\ref{sec:toolkit-odds}. In words: the weight that a top posterior leaves on class \(L\) equals the bottom
posterior of a market whose prior on class \(H\) is \(1-\pi\). The two classes swap roles.

\begin{derivation}[fact 2: \(m\) rises strictly]
For a prior \(x\in(0,1)\),
\begin{align*}
m(x) &= \frac{x}{x+(1-x)(1+u)}=\frac{x}{x+1+u-x-xu} \why{expand \((1-x)(1+u)\)}\\
&= \frac{x}{1+u(1-x)} \why{\(x-x=0\); \(u-xu=u(1-x)\)}
\end{align*}
As \(x\) rises, the numerator rises and the denominator falls, because \(u>0\). Both stay positive, so \(m(x)\)
rises strictly.
\end{derivation}

With the two facts,
\begin{align*}
\kappa(\pi) &= \min\{m(\pi),\,m(1-\pi)\} \why{fact 1}\\
&= m\bigl(\min\{\pi,1-\pi\}\bigr) \why{fact 2: the smaller argument gives the smaller value}\\
&\le m\bigl(\tfrac12\bigr) \why{\(\min\{\pi,1-\pi\}\le\tfrac12\); fact 2}\\
&= \frac{1/2}{1+u/2}=\frac{1}{2+u} \why{fact 2 form at \(x=\tfrac12\); multiply above and below by 2}
\end{align*}
The last inequality is an equality only at \(\pi=\tfrac12\), because \(m\) rises strictly. So the bound is largest
at the prior \(\tfrac12\), where it equals \(1/(2+u)\), and smaller at every other prior
(Figure~\ref{fig:trading-kappa}).

\begin{figure}[htbp]
\centering
\begin{tikzpicture}[x=9cm,y=4.4cm]
\draw[-{Stealth[length=2mm]}] (0,0) -- (1.07,0) node[right] {\(\pi\)};
\draw[-{Stealth[length=2mm]}] (0,0) -- (0,1.08);
\node[below] at (0,0) {\(0\)};
\draw (1,0.02) -- (1,-0.02) node[below] {\(1\)};
\draw (0.5,0.02) -- (0.5,-0.02) node[below] {\(\tfrac12\)};
\draw (0.01,1) -- (-0.01,1) node[left] {\(1\)};
\draw[rulegrey, domain=0.001:0.999, samples=80, smooth] plot (\x, {\x/(1+1.5*(1-\x))});
\draw[rulegrey, domain=0.001:0.999, samples=80, smooth] plot (\x, {(1-\x)/(1+1.5*\x)});
\node[right, notegrey] at (1,0.8) {\(m(\pi)\)};
\node[left, notegrey] at (0,0.8) {\(1-M(\pi)\)};
\draw[accent, very thick, domain=0.001:0.5, samples=60, smooth] plot (\x, {\x/(1+1.5*(1-\x))});
\draw[accent, very thick, domain=0.5:0.999, samples=60, smooth] plot (\x, {(1-\x)/(1+1.5*\x)});
\draw[dotted] (0.5,0) -- (0.5,0.2857);
\draw[dotted] (0,0.2857) -- (0.5,0.2857);
\node[left] at (0,0.2857) {\(\tfrac{1}{2+u}\)};
\node[accent, above] at (0.5,0.29) {\(\kappa(\pi)\)};
\end{tikzpicture}
\caption{The bound \(\kappa(\pi)\) (thick line) is the lower of the grey curves \(m(\pi)\) and \(1-M(\pi)\). Left
of \(\tfrac12\), class \(L\)'s floor binds; right of \(\tfrac12\), class \(H\)'s floor binds. The peak is at the
prior \(\tfrac12\). The curves are drawn for one value of \(u\); the shape is the same for every \(u>0\).}
\label{fig:trading-kappa}
\end{figure}

\paragraph{Which class binds.} If the prior favors class \(H\), \(\pi>\tfrac12\), then \(1-\pi<\pi\) and
\(\kappa(\pi)=m(1-\pi)=1-M(\pi)\), the floor of \(K_H\). If \(\pi<\tfrac12\), then \(\kappa(\pi)=m(\pi)\), the floor
of \(K_L\). Either way the favored class binds, because the price already leans toward it. When \(\pi\) is high, a
purchase pushes a posterior that is already high even closer to one. Little weight on class \(L\) remains, so little
is left for class \(H\) to earn. A sale by class \(L\) starts from the same high prior and has much more to correct.

\paragraph{Extreme cases.} When the investor is rare relative to liquidity (\(u\) near zero), \(m(\pi)\) and
\(M(\pi)\) are near \(\pi\), so \(\kappa(\pi)\) is near \(\min\{\pi,1-\pi\}\). When the investor is common (\(u\)
large), a purchase or a sale can reveal almost all that the investor knows: \(m(\pi)\) is near zero and \(M(\pi)\)
is near one, so \(\kappa(\pi)\) is near zero. When the prior is near zero or one,
\(\kappa(\pi)\) is near zero, because one class is almost certain before any trade.

\begin{checkyourself}\label{q:trading-prior}
Let \(\pi>\tfrac12\). Show, with each step, that \(1-M(\pi)<m(\pi)\). Then say in words which investor has the
smaller guaranteed return, and why.
\end{checkyourself}

\subsection{What a bound for every candidate buys}

The bound holds under every candidate profile, mixed or pure. Its derivation used only three facts: entry lies in
\([\rho,1]\), the spread is nonnegative, and the posteriors after nonzero orders lie in \([m(\pi),M(\pi)]\). None of
them uses the investor's optimality (E5) or a guess about the profile. So the bound constrains every candidate
before anyone solves for the equilibrium. This is how the paper proves uniqueness: a uniqueness proof derives its
bounds for every candidate and does not use the proposed tuple. Chapter~\ref{ch:ccc} uses both sides.
\begin{itemize}
\item If \(\DT<k\), then \(K_j\le\DT<k\) in every candidate. By parts (a) and (b), each class abstains, so no
equilibrium has informed trade.
\item If \(\rho\,\kappa(\pi)\,\DT>k\), then \(K_j>k\) in every candidate. By parts (a) and (b), class \(H\) buys for
sure and class \(L\) sells for sure in every equilibrium.
\end{itemize}
When \(\rho\,\kappa(\pi)\,\DT\le k\le\DT\), the bound alone does not decide the orders. Chapter~\ref{ch:coexistence}
studies that region.

\begin{watchout}
Keep the domain of \(\beta\) of each result.
\begin{itemize}
\item The trading bound uses the posterior bounds only after the nonzero orders (Lemma A.1(b)), so it holds for
every \(0<\beta<1/2\).
\item The restriction \(\beta\le1/3\) is the inequality \(b_0\ge b\), where \(b_0=(1-\lambda)(1-2\beta)\) is the
probability of a liquidity abstention. It extends the posterior bounds to the zero order, and through this it enters
only the exclusion bound (Corollary A.1) and the zero-order part of the research bounds. The research parameter
region and the slot-feasibility parts of identification also range over the slots with \(\beta\le1/3\).
\item A price reached by at least one nonzero order gives a belief in \([m(\pi),M(\pi)]\) for every
\(0<\beta<1/2\), even when the price pools that order with the zero order (Section~\ref{sec:beliefs-price-belief}).
\end{itemize}
The lower bound is also positive only because \(\rho>0\): with \(\rho=0\), a purchase that does not recruit would
earn nothing.
\end{watchout}

% =====================================================================
\section{The problem splits by class}\label{sec:trading-split}

Part (d) of Lemma A.3 is visible in \eqref{eq:trading-returns}. \(K_H\) contains only \(\mu(1)\) and \(e(1)\), and
\(K_L\) only \(\mu(-1)\) and \(e(-1)\). Entry after the zero order enters no payoff: abstention earns zero whatever
happens after the zero order, because \(q=0\) multiplies the bracket by zero.

Combine part (d) with part (a). In equilibrium, class \(L\) never buys and class \(H\) never sells. So
\(\mu(1)\) depends only on how often class \(H\) buys (Section~\ref{sec:trading-best}), and \(\mu(-1)\) only on how
often class \(L\) sells. When \(\DT>0\) and \(g_H>g_L\), entry after each order follows its own posterior through
the threshold rule. So:
\begin{itemize}
\item the conditions of class \(H\) involve only its purchase rate, \(\mu(1)\) and \(e(1)\);
\item the conditions of class \(L\) involve only its sale rate, \(\mu(-1)\) and \(e(-1)\).
\end{itemize}
The investor's problem splits into two separate problems, one for each class. The classes still meet at the zero
order, whose posterior depends on both rates, but nothing after the zero order enters any payoff.
Chapter~\ref{ch:coexistence} uses this split: it solves the purchase side and the sale side one at a time.

% =====================================================================
\section{The two numbers}\label{sec:trading-twonumbers}

Look back at what the investor's problem uses. The payoffs use the auction only through the spread \(\DT\). The
entry rule uses it only through the threshold. Recall \eqref{eq:beliefs-threshold}: when \(g_H>g_L\), the
threshold \(\tau=(c_H-g_L)/(g_H-g_L)\) is the belief at which gross acquisition profit
\(B(\mu)=g_L+\mu\,(g_H-g_L)\) equals the high preparation cost \(c_H\). So the spread sets what an informed order
earns at each posterior and level of entry, and the threshold sets which posteriors recruit a high-cost
challenger. Write \(\gamma=g_H-g_L\) for the profit gap between the classes.

\begin{paperresult}{Lemma A.4 (The auction through the spread and the threshold)}
Fix the model's independence and information structure, a prior \(\pi\in(0,1)\), the two-point preparation cost
with \(\rho\in(0,1)\) and \(c_H>0\), a trading cost \(k>0\), and one pair of tie rules: an indifferent investor may
mix over its best orders or abstains, and an indifferent high-cost challenger prepares or may prepare with any
probability. At fixed presence, fix \(0<\lambda<1\); with research, presence is \(\hat\lambda\). Fix
\(0<\beta<1/2\). Consider two cash second-price auctions with finite-mean classes satisfying the class ranking,
continuous incumbent laws, \(\DT>0\) and \(\gamma=g_H-g_L>0\), and the same spread and preparation threshold
\(\tau=(c_H-g_L)/\gamma\). Their equilibrium sets projected onto \((\sigma_H,\sigma_L,\mu,e)\) coincide.

The common conditions are Bayes' rule, the threshold entry rule under the chosen tie rule, and optimal orders at
\[
\begin{aligned}
U_H(1)&=e(1)\DT[1-\mu(1)]-k,& U_H(-1)&=-e(-1)\DT[1-\mu(-1)]-k,\\
U_L(-1)&=e(-1)\DT\mu(-1)-k,& U_L(1)&=-e(1)\DT\mu(1)-k,
\end{aligned}
\qquad U_j(0)=0.
\]
Each auction gives prices and price beliefs by
\[
P(x)=t_0+e(x)[w_L+\DT\mu(x)],\qquad \zeta(P(x))=\frac{e(x)-\rho}{1-\rho},\qquad \bar\mu(P(x))=\mu(x).
\]
With research, fix also the quadratic cost coefficient \(a>0\) and cap \(0<\lb<1\), and require optimal orders
after a hypothetical success even at zero research. The equilibrium sets projected onto
\((\hat\lambda,\sigma_H,\sigma_L,\mu,e)\) coincide. The common conditions also include
\[
\hat\lambda=\min\{\lb,\Psi/a\},\qquad \Psi=\pi\max_q U_H(q)+(1-\pi)\max_q U_L(q),
\]
with Bayes' rule evaluated at \(\hat\lambda\). These conclusions hold under all four tie rules.
\end{paperresult}

A plain reading:
\begin{itemize}
\item \emph{Held fixed:} everything outside the auction, that is, the information structure, the prior, \(\rho\)
and \(c_H\), the trading cost \(k\), the slot, and one choice of tie rules for the investor and the challenger.
The guide reads the phrase ``every other primitive and the tie rules the same'' of Section 3.4 of the paper as
this list.
\item \emph{Not on the fixed list:} the classes, the reserve and the incumbent's law. The statement asks only that
each auction have finite-mean classes with the class ranking, a continuous incumbent law, \(\DT>0\) and
\(\gamma>0\), and that the two auctions share \(\DT\) and \(\tau\). Their levels \(t_0\), \(t_L\), \(g_L\) and
\(g_H\) can differ.
\item \emph{The conclusion:} drop from each equilibrium its prices, price beliefs and preparation rule, and keep
the orders, the posteriors and entry. The two auctions then have exactly the same list of such objects. These
objects solve the common conditions, which contain \(\DT\) and \(\tau\) and no other feature of the auction. Each
auction builds its own prices from its own \(t_0\) and \(w_L=t_L-t_0\).
\item \emph{Research:} when the investor buys its presence \(\hat\lambda\) at a quadratic cost, the same holds with
\(\hat\lambda\) added to the list.
\end{itemize}

\subsection{The proof}

The proof has two directions: every equilibrium of either auction solves the common conditions, and every solution
of the common conditions is part of an equilibrium in each auction.

\paragraph{Step 1: the price reveals the posterior.} In either auction \(\DT>0\), \(\gamma>0\) and \(\rho>0\), so
Lemma A.2(c) applies (Section~\ref{sec:beliefs-reveal}). The challenger's belief at the price \(P(x)\) equals
\(\mu(x)\).

\paragraph{Step 2: the entry rule uses only \(\tau\).}
\begin{align*}
B(\mu)-c_H &= g_L+\mu\,\gamma-c_H \why{\(B(\mu)=g_L+\mu\gamma\)}\\
&= \gamma\,\Bigl[\mu-\frac{c_H-g_L}{\gamma}\Bigr]=\gamma\,(\mu-\tau) \why{factor out \(\gamma>0\); definition of \(\tau\)}
\end{align*}
Since \(\gamma>0\), \(B(\mu)-c_H\) has the sign of \(\mu-\tau\). With Step 1, entry optimality (E3) reads:
\begin{itemize}
\item if \(\mu(x)<\tau\), then \(e(x)=\rho\);
\item if \(\mu(x)>\tau\), then \(e(x)=\rho+(1-\rho)=1\);
\item if \(\mu(x)=\tau\), then \(e(x)=1\) when an indifferent high-cost challenger prepares, and \(e(x)\) can take
any value in \([\rho,1]\) when it may prepare with any probability.
\end{itemize}
This threshold rule uses \(\tau\), \(\rho\) and the tie rule, and nothing else from the auction.

\paragraph{Step 3: the payoffs use only \(\DT\).} Lemma A.3 gives the four nonzero-order payoffs, and
\(U_j(0)=0\). The levels \(t_0\) and \(w_L\) cancel (Section~\ref{sec:trading-returns}).

\paragraph{Step 4: Bayes' rule does not use the auction.} Posteriors come from the prior, the slot and the profile.

So every equilibrium of either auction gives orders, posteriors and entry that solve the common conditions, and
these conditions are the same in both auctions. This is the first direction.

\paragraph{Step 5: from a solution back to an equilibrium.} Take orders, posteriors and entry that solve the common
conditions, and fix one auction. Build its prices, price beliefs and preparation rule, and check (E1) to (E6).

\emph{5a. Entry does not fall when the posterior rises.} This is the entry step of Lemma A.2(e)
(Section~\ref{sec:beliefs-reveal}), applied to the threshold rule of Step 2. Let \(\mu(x)>\mu(x')\). If \(\mu(x')<\tau\), then
\(e(x')=\rho\le e(x)\). If \(\mu(x')\ge\tau\), then \(\mu(x)>\tau\), so \(e(x)=1\ge e(x')\).

\emph{5b. Prices, preparation rule and price beliefs.} Set \(P(x)=t_0+e(x)\,[w_L+\DT\mu(x)]\) with this
auction's \(t_0\) and \(w_L\). Then set \(\zeta(P(x))=[e(x)-\rho]/(1-\rho)\) and \(\bar\mu(P(x))=\mu(x)\). The
claim on constructed tuples (Section~\ref{sec:beliefs-constructed}) needs \(\DT>0\), entry in \([\rho,1]\), and
entry that does not fall when the posterior rises. Step 5a gives the last of these. The claim then gives:
\begin{itemize}
\item two orders share a price exactly when they share the posterior and entry, by the increment identity
\eqref{eq:beliefs-increment};
\item so \(\zeta\) is a function of the price, with values in \([0,1]\);
\item so \(\bar\mu\) is well defined, and it equals the posterior of the event \(\{x:P(x)=\bar P\}\), because all
orders in this event share one likelihood ratio.
\end{itemize}
Under the alternative challenger rule, two orders at \(\tau\) can have different entry. They then get different
prices, and each of these prices reveals the posterior \(\tau\).

\emph{5c. The six conditions.} Check each condition by direct substitution.
\begin{itemize}
\item (E1) is a common condition.
\item (E2) is the last point of 5b.
\item (E3) holds by Step 2, since the belief at each price is \(\mu(x)\).
\item (E4) holds because Steps 1 to 4 of Section~\ref{sec:beliefs-priceformula} compute
\(\E[V_T\mid X=x]=t_0+e(x)\,w(\mu(x))\) from (E1), the truthful bids and the independence of the cost. This
computation does not use (E4), and the price of 5b equals its result.
\item (E5) holds because, by Lemma A.3, the payoffs at these prices are the payoffs of the common conditions.
\item (E6) holds because truthful bids are weakly dominant (Chapter~\ref{ch:auction}).
\end{itemize}
So the tuple is an equilibrium of this auction. The same construction works in the other auction, so the projected
equilibrium sets coincide.

\paragraph{Step 6: research.} The investor chooses an intensity \(\tilde\lambda\in[0,\lb]\) at the cost
\((a/2)\tilde\lambda^2\) before it learns the class. Market makers and the challenger do not see this choice, so
prices and entry stay fixed. Research succeeds with probability \(\tilde\lambda\); the investor then learns the
class (\(H\) with probability \(\pi\)) and places its best order. After a failure a liquidity trader fills the slot
and the investor earns zero. The expected payoff is
\[
\tilde\lambda\bigl[\pi\max_qU_H(q)+(1-\pi)\max_qU_L(q)\bigr]-\frac a2\tilde\lambda^2=\tilde\lambda\,\Psi-\frac a2\tilde\lambda^2 .
\]
Its derivative \(\Psi-a\tilde\lambda\) falls strictly, so the payoff is strictly concave. Also \(\Psi\ge0\), because
\(U_j(0)=0\) gives \(\max_qU_j(q)\ge0\). If \(\Psi/a\le\lb\), the first-order condition \(\Psi-a\tilde\lambda=0\)
gives \(\tilde\lambda=\Psi/a\in[0,\lb]\). If \(\Psi/a>\lb\), the derivative is positive on all of \([0,\lb]\), and
the best choice is \(\lb\). Together, \(\hat\lambda=\min\{\lb,\Psi/a\}\). The value \(\Psi\) uses only the payoffs of
Lemma A.3, and \(a\) and \(\lb\) are common, so the research condition is the same in both auctions. At zero
intensity, the required optimality after a hypothetical success gives the same best orders. Bayes' rule uses
\(\hat\lambda\) as the presence probability in both.

\begin{watchout}
Keep the hypotheses of Lemma A.4.
\begin{itemize}
\item \emph{A positive spread.} Steps 1 and 5b need \(\DT>0\), with \(\rho>0\). Otherwise the price need not reveal
the posterior.
\item \emph{A positive profit gap.} Step 2 needs \(\gamma>0\). With \(\gamma=0\), gross profit does not depend on the
belief, and no threshold exists.
\item \emph{The same tie rules and other primitives.} Entry at \(\mu=\tau\) and the choice at \(K_j=k\) depend on
the tie rules. The lemma holds under each of the four pairs, but both auctions must use the same pair.
\end{itemize}
\end{watchout}

\subsection{What the lemma means}

Orders, posteriors, entry and research depend on the auction only through two numbers, the spread \(\DT\) and the
threshold \(\tau\). Two other objects need more.
\begin{itemize}
\item \emph{Prices} add the levels \(t_0\) and \(t_L\): \(P(x)=t_0+e(x)\,[t_L-t_0+\DT\,\mu(x)]\).
\item \emph{Ownership} adds who wins. High-class ownership is the prior times entry in class \(H\) times class
\(H\)'s chance to win against the incumbent, \(W_H(F)\). So every result on ownership carries \(W_H(F)\). This is why
the clean rise of high-class ownership in Chapter~\ref{ch:ccc} needs class \(H\) to outbid the stronger incumbent,
\(W_H(F_1)=1\).
\end{itemize}

\begin{checkyourself}\label{q:trading-levels-two}
Two auctions satisfy the hypotheses of Lemma A.4 with the same \(\DT\) and \(\tau\), but different \(t_0\) and
\(t_L\). Which equilibrium objects agree, and which can differ? Write the difference of the two prices after an
order \(x\) in symbols.
\end{checkyourself}

% =====================================================================
\section{The two numbers along a path}\label{sec:trading-path}

Recall that a path is a family \(\{F_r\}\) of continuous incumbent laws, indexed by a strength \(r\) in an interval.
Whenever \(r'>r\), the law \(F_{r'}\) is a first-order stochastic strengthening of \(F_r\)
(Section~\ref{sec:toolkit-fosd}): \(F_{r'}(y)\le F_r(y)\) at every \(y\).

\begin{paperresult}{Lemma A.5 (The two numbers along a path)}
Fix ranked finite-mean classes, the reserve and the two-point preparation cost. A path is a family \(F_r\) of
continuous incumbent laws indexed by strength on an interval, with \(F_{r'}\le F_r\) when \(r'>r\). Write
\(\DT(r)=\DT(F_r)\), \(\gamma(r)=g_H(F_r)-g_L(F_r)\), and, where \(\gamma(r)>0\),
\(\tau(r)=[c_H-g_L(F_r)]/\gamma(r)\). The class mean gap is fixed:
\[
\bar\Gamma=\E_H\theta-\E_L\theta=\DT(r)+\gamma(r).
\]
Each \(g_j\) and \(\gamma\) is nonincreasing, and \(\DT\) is nondecreasing. For \(r'>r\),
\[
\begin{aligned}
\DT(r')-\DT(r)&=\int_p^\infty[F_r(y)-F_{r'}(y)]D(y)\,dy,\\
g_j(F_r)-g_j(F_{r'})&=\int_p^\infty[F_r(y)-F_{r'}(y)][1-G_j(y)]\,dy.
\end{aligned}
\]
Each change is strict exactly when its integral is positive. Where \(\gamma>0\), \(\DT>0\) and \(\tau\ge0\), both
\(\tau\) and \(\DT\tau\) are nondecreasing. Their exact increments are
\[
\tau(r')-\tau(r)=\frac{g_L(F_r)-g_L(F_{r'})+\tau(r)[\DT(r')-\DT(r)]}{\gamma(r')},
\]
\[
\DT(r')\tau(r')-\DT(r)\tau(r)=\DT(r)[\tau(r')-\tau(r)]+\tau(r')[\DT(r')-\DT(r)].
\]
Thus the threshold rises strictly exactly when the numerator in its increment is positive; the product rises
strictly exactly when the right side of its increment is positive.

On a differentiable path the return \(J=\DT(1-\tau)\) obeys
\[
J'=\frac{\bar\Gamma(g_H-c_H)\DT'+\DT\gamma g_H'}{\gamma^2}.
\]
Its direction depends on the balance between spread growth and the decline in high-class profit. The paper gives an
instance with point classes \(h\) and \(\ell\), a reserve \(p\), a high cost \(c_H\), and \(F_r\) uniform on a short
interval centered at \(r\). The strengths \(r\) run over an interval on which every incumbent value lies strictly
between \(\ell\) and \(h\). There
\[
\DT=r-\ell,\qquad g_L=0,\qquad g_H=h-r,\qquad \tau=\frac{c_H}{h-r},
\]
and the return \(J\) is lower at the right end of the interval of strengths than at the left end. In the
uniform-incumbent example, \(J\) rises on \([r_0,r_2]\) by the uniform case of Proposition~3(c).
\end{paperresult}

The paper states the instance with particular values of \(h\), \(\ell\), \(p\), \(c_H\), the interval and the two end
values of \(J\). The box above gives it in words. Chapter~\ref{app:numbers} lists the values; you do not need them.

A plain reading:
\begin{itemize}
\item Only the incumbent changes along the path, and the sum of the spread and the profit gap stays fixed.
\item A stronger incumbent never raises either class's gross profit, never widens the profit gap, and never narrows
the spread. The two integrals measure each change exactly.
\item Where the profit gap and the spread are positive and the threshold is not negative, a stronger incumbent
never lowers the threshold, and the product \(\DT\tau\) never falls.
\item The number \(J=\DT(1-\tau)\) can rise or fall. The instance shows a fall; the paper's example shows a rise.
\end{itemize}

What is \(J\)? Read it off \eqref{eq:trading-returns}. A purchase whose posterior just reaches the threshold,
\(\mu(1)=\tau\), recruits a high-cost challenger under the default tie rules, so \(e(1)=1\) and
\(K_H=\DT\,(1-\tau)=J\). So \(J\) is the gross return of a purchase that just reaches \(\tau\) and recruits.
Chapter~\ref{ch:coexistence} uses \(J\), and also the product \(\DT\tau\).

\paragraph{Step 1: the class mean gap is fixed.} Recall the class-gap identity \eqref{eq:auction-gap}. Under the
class ranking every value meets the reserve, so \(\max\{\theta,p\}=\theta\) with probability one, and
\begin{align*}
\DT(r)+\gamma(r) &= \E_H[\max\{\theta,p\}]-\E_L[\max\{\theta,p\}] \why{class-gap identity at \(F_r\)}\\
&= \E_H\theta-\E_L\theta=\bar\Gamma \why{every value meets the reserve}
\end{align*}
The right side does not contain \(F_r\). The paper's proof reaches the same point through the split
\(t_j+g_j=\E_j\theta\) of Chapter~\ref{ch:auction}: subtract the split of class \(L\) from that of class \(H\). A
stronger incumbent moves the class value gap between the profit gap and the spread and leaves the total unchanged.

\paragraph{Step 2: the profit changes.} Recall the gross profit integral \eqref{eq:auction-gross},
\(g_j(F)=\int_p^\infty F(y)\,[1-G_j(y)]\,dy\). It adds up, over levels \(y\) above the reserve, the probability that
the incumbent's value lies below \(y\) times the probability that the challenger's value lies above \(y\). Then
\begin{align*}
\gamma(F) &= \int_p^\infty F(y)\,\bigl[G_L(y)-G_H(y)\bigr]\,dy=\int_p^\infty F(y)\,D(y)\,dy \why{\(g_H-g_L\); the 1s cancel}\\
g_j(F_r)-g_j(F_{r'}) &= \int_p^\infty\bigl[F_r(y)-F_{r'}(y)\bigr]\,[1-G_j(y)]\,dy\ge0 \why{\(F_r\ge F_{r'}\), \(1-G_j\ge0\)}
\end{align*}
This is Proposition~1(a). In the same way, \(\gamma(r)-\gamma(r')=\int_p^\infty[F_r-F_{r'}]\,D\,dy\ge0\), because
\(D\ge0\) under the class ranking.

\paragraph{Step 3: the spread change.}
\begin{align*}
\DT(r')-\DT(r) &= \bigl[\bar\Gamma-\gamma(r')\bigr]-\bigl[\bar\Gamma-\gamma(r)\bigr]=\gamma(r)-\gamma(r') \why{Step 1 at \(r'\) and \(r\)}\\
&= \int_p^\infty\bigl[F_r(y)-F_{r'}(y)\bigr]\,D(y)\,dy\ge0 \why{Step 2; \(F_r\ge F_{r'}\), \(D\ge0\)}
\end{align*}
This is Proposition~1(b). Each integrand in Steps 2 and 3 is nonnegative, so each change is zero or positive, and
strictly positive exactly when its integral is positive. Proposition~1 gives weak inequalities; strictness needs
these integral conditions.

\paragraph{Step 4: where the threshold is defined and nonnegative.} The threshold is defined where \(\gamma(r)>0\).
There,
\[
\tau(r)\ge0 \iff c_H-g_L(F_r)\ge0 \iff g_L(F_r)\le c_H ,
\]
because multiplication by \(\gamma(r)>0\) keeps the direction. So the lemma's condition \(\tau\ge0\) is the
condition \(g_L\le c_H\): a challenger known to be in class \(L\) earns at most the high preparation cost. The two
conditions persist in opposite directions along a path.
\begin{itemize}
\item \(g_L\le c_H\) persists toward stronger incumbents: \(g_L\) is nonincreasing (Step 2), so
\(g_L(F_{r'})\le g_L(F_r)\le c_H\) for \(r'>r\). Once it holds, it holds against every stronger incumbent.
\item \(\gamma>0\) persists toward weaker incumbents: \(\gamma\) is nonincreasing, so \(\gamma(r')>0\) gives
\(\gamma(r)\ge\gamma(r')>0\) for \(r<r'\). The threshold need not be defined along the whole path.
\end{itemize}
So to compare \(r<r'\), the lemma needs \(\gamma(r')>0\) and \(\tau(r)\ge0\).

\paragraph{Step 5: the threshold never falls.} First a picture. Gross profit \(B_r(\mu)=g_L(F_r)+\mu\,\gamma(r)\) is
a straight line in the belief: height \(g_L\) at \(\mu=0\), height \(g_H\) at \(\mu=1\), slope \(\gamma\). The
threshold is where it crosses \(c_H\). A stronger incumbent lowers the intercept and flattens the line, so the line
reaches \(c_H\) further right (Figure~\ref{fig:trading-threshold}).

\begin{figure}[htbp]
\centering
\begin{tikzpicture}[x=10cm,y=0.95cm]
\draw[-{Stealth[length=2mm]}] (0,0) -- (1.1,0) node[right] {\(\mu\)};
\draw[-{Stealth[length=2mm]}] (0,0) -- (0,4.6);
\node[below] at (0,0) {\(0\)};
\draw (1,0.06) -- (1,-0.06) node[below] {\(1\)};
\draw[dashed, notegrey] (0,2.4) -- (1.02,2.4) node[right] {\(c_H\)};
\draw[accent, thick] (0,1.0) -- (1,4.2) node[right] {\(B_r\)};
\draw[ink, thick] (0,0.4) -- (1,3.2) node[right] {\(B_{r'}\)};
\node[left] at (0,1.0) {\(g_L(F_r)\)};
\node[left] at (0,0.4) {\(g_L(F_{r'})\)};
\draw[dotted] (0.4375,2.4) -- (0.4375,0) node[below] {\(\tau(r)\)};
\draw[dotted] (0.7143,2.4) -- (0.7143,0) node[below] {\(\tau(r')\)};
\end{tikzpicture}
\caption{Gross profit as a function of the belief, against a weaker incumbent (\(B_r\)) and a stronger one
(\(B_{r'}\)). The stronger incumbent lowers the intercept \(g_L\) and the slope \(\gamma\). When the threshold \(\tau(r)\)
is not negative, the line \(B_{r'}\) at that belief lies at or below \(c_H\), so it meets \(c_H\) further right. The figure is the guide's own, and the positions are only a picture.}
\label{fig:trading-threshold}
\end{figure}

The guide's own short argument shows the condition hidden in the picture. At the threshold of the weaker
incumbent, \(B_r(\tau(r))=c_H\), and
\begin{align*}
B_{r'}(\tau(r)) &= g_L(F_{r'})+\tau(r)\,\gamma(r') \why{definition of \(B_{r'}\)}\\
&\le g_L(F_r)+\tau(r)\,\gamma(r') \why{\(g_L(F_{r'})\le g_L(F_r)\)}\\
&\le g_L(F_r)+\tau(r)\,\gamma(r) \why{\(\gamma(r')\le\gamma(r)\) and \(\tau(r)\ge0\)}\\
&= c_H \why{definition of \(\tau(r)\)}
\end{align*}
The line \(B_{r'}\) rises with \(\mu\) and is at or below \(c_H\) at \(\tau(r)\), so it reaches \(c_H\) at a belief at
or above \(\tau(r)\). The third line needs \(\tau(r)\ge0\): with a negative \(\tau(r)\), the product of \(\tau(r)\)
and the smaller slope is the \emph{larger} number.

The paper's proof gives the exact increment, with \(\gamma(r')=\gamma(r)-[\DT(r')-\DT(r)]\) from Steps 1 and 3.

\begin{derivation}[the increment of the threshold]
\begin{align*}
\tau(r)\,\gamma(r') &= \tau(r)\,\gamma(r)-\tau(r)\,[\DT(r')-\DT(r)] \why{Steps 1 and 3}\\
&= c_H-g_L(F_r)-\tau(r)\,[\DT(r')-\DT(r)] \why{\(\tau(r)\gamma(r)=c_H-g_L(F_r)\)}\\[4pt]
\tau(r')-\tau(r) &= \frac{c_H-g_L(F_{r'})-\tau(r)\,\gamma(r')}{\gamma(r')} \why{common denominator \(\gamma(r')\)}\\
&= \frac{c_H-g_L(F_{r'})-c_H+g_L(F_r)+\tau(r)\,[\DT(r')-\DT(r)]}{\gamma(r')} \why{line 2}\\
&= \frac{g_L(F_r)-g_L(F_{r'})+\tau(r)\,[\DT(r')-\DT(r)]}{\gamma(r')} \why{\(c_H-c_H=0\)}
\end{align*}
\end{derivation}

In the numerator, \(g_L(F_r)-g_L(F_{r'})\ge0\) by Proposition~1(a), \(\tau(r)\ge0\) by assumption, and
\(\DT(r')-\DT(r)\ge0\) by Proposition~1(b). The denominator \(\gamma(r')\) is positive. So \(\tau(r')\ge\tau(r)\).
The threshold rises strictly exactly when the numerator is positive: when \(g_L\) falls strictly, or when the
spread widens strictly at a positive threshold.

\begin{intuition}[why a stronger incumbent raises the threshold]
A stronger incumbent never raises class \(L\)'s profit and never widens the profit gap between the classes. So a
challenger needs at least as favorable a belief before preparation at the high cost pays. The profit gap shrinks
because the class value gap moves from the challenger's profit to target proceeds (Step 1).
\end{intuition}

\paragraph{Step 6: why the condition \(g_L\le c_H\) is needed.} Suppose \(g_L(F_r)>c_H\). Then \(\tau(r)<0\), and
\(B_r(\mu)\ge g_L(F_r)>c_H\) at every belief, so a high-cost challenger prepares after every price and entry is
one after every order. The threshold does not bind. In this region the second term of the increment,
\(\tau(r)\,[\DT(r')-\DT(r)]\), is negative when the spread widens strictly. If \(g_L\) does not change, that term
is the whole numerator, and the threshold falls. So ``a stronger incumbent never lowers the threshold'' holds only
while \(g_L\le c_H\).

\begin{watchout}
Keep two conditions with the statement on the threshold.
\begin{itemize}
\item \(g_L\le c_H\), that is \(\tau\ge0\). Without it the threshold can fall. Once it holds on a path, it holds at
every stronger incumbent, because \(g_L\) never rises.
\item \(\gamma>0\), so that the threshold exists. This condition can fail as the incumbent grows stronger, because
\(\gamma\) never rises. As \(\gamma\) falls toward zero with \(c_H-g_L>0\), the threshold grows without bound and
soon exceeds one; then no posterior recruits.
\end{itemize}
\end{watchout}

\paragraph{Step 7: the product \(\DT\tau\).}
\begin{align*}
&\DT(r')\tau(r')-\DT(r)\tau(r)\\
&\quad= \DT(r')\tau(r')-\DT(r)\tau(r')+\DT(r)\tau(r')-\DT(r)\tau(r) \why{add and subtract \(\DT(r)\tau(r')\)}\\
&\quad= \tau(r')\,[\DT(r')-\DT(r)]+\DT(r)\,[\tau(r')-\tau(r)] \why{factor each pair}
\end{align*}
Both terms are nonnegative, since \(\tau(r')\ge\tau(r)\ge0\), \(\DT(r')\ge\DT(r)>0\) and \(\tau(r')-\tau(r)\ge0\).
So the product never falls, and it rises strictly exactly when the sum is positive.

\paragraph{Step 8: the return of a purchase that just recruits.} \(J=\DT(1-\tau)\) has one factor that never falls
and one that never rises, so its direction is not fixed. First,
\[
1-\tau=\frac{\gamma-(c_H-g_L)}{\gamma}=\frac{g_H-g_L-c_H+g_L}{\gamma}=\frac{g_H-c_H}{\gamma},
\]
so \(J=\DT\,(g_H-c_H)/\gamma\). Along the path \(\gamma=\bar\Gamma-\DT\) with \(\bar\Gamma\) fixed, so
\(\gamma'=-\DT'\).

\begin{derivation}[the derivative of \(J\)]
\begin{align*}
J' &= \frac{(g_H-c_H)\DT'+\DT g_H'}{\gamma}+\DT(g_H-c_H)\Bigl(-\frac{\gamma'}{\gamma^2}\Bigr) \why{product rule; \(c_H\) fixed}\\
&= \frac{(g_H-c_H)\DT'+\DT g_H'}{\gamma}+\frac{\DT(g_H-c_H)\DT'}{\gamma^2} \why{\(\gamma'=-\DT'\)}\\
&= \frac{(g_H-c_H)\DT'\gamma+\DT\gamma g_H'+\DT(g_H-c_H)\DT'}{\gamma^2} \why{common denominator \(\gamma^2\)}\\
&= \frac{(g_H-c_H)\DT'(\gamma+\DT)+\DT\gamma g_H'}{\gamma^2} \why{collect the \(\DT'\) terms}\\
&= \frac{\bar\Gamma(g_H-c_H)\DT'+\DT\gamma g_H'}{\gamma^2} \why{\(\gamma+\DT=\bar\Gamma\)}
\end{align*}
\end{derivation}

Where \(g_H>c_H\), the first term is the gain from spread growth (\(\DT'\ge0\)), and it is nonnegative. The second
term is the loss from the decline in high-class profit (\(g_H'\le0\)), and it is nonpositive. The direction of
\(J\) depends on which is larger.

\paragraph{Step 9: the instance in symbols.} Take point classes: the value is \(h\) in class \(H\) and \(\ell\) in
class \(L\), with \(p<\ell<h\). Let the incumbent be uniform on a short interval centered at \(r\), with every
incumbent value strictly between \(\ell\) and \(h\).
\begin{itemize}
\item Class \(L\): the incumbent always wins and pays the challenger's bid \(\ell\), which meets the reserve. So
\(t_L=\ell\) and \(g_L=0\).
\item Class \(H\): the challenger always wins and pays \(\max\{p,R\}=R\), because \(R>\ell>p\). So \(t_H=\E R=r\)
and \(g_H=\E[h-R]=h-r\).
\end{itemize}
Hence \(\DT=r-\ell\), \(\gamma=h-r\) and \(\tau=c_H/(h-r)\). As a check, \(\DT+\gamma=h-\ell=\bar\Gamma\), free of
\(r\).

The guide's own reading of these formulas shows the shape of \(J=(r-\ell)\bigl(1-c_H/(h-r)\bigr)\) without numbers.
It holds only where the formulas hold, that is, while every incumbent value stays strictly between \(\ell\) and
\(h\). The first factor rises with \(r\). The second falls with \(r\) and reaches zero at \(r=h-c_H\). So, if the
path reaches near \(h-c_H\) with the formulas still valid, \(J\) falls toward zero there, even though the spread
continues to rise. The derivative agrees: with \(\DT'=1\), \(g_H'=-1\) and \(\bar\Gamma=h-\ell\), the numerator of
\(J'\) is
\[
(h-\ell)(h-r-c_H)-(r-\ell)(h-r).
\]
Near \(r=h-c_H\) the first product is near zero and the second is positive, so the numerator is negative. The paper
picks particular values and finds \(J\) lower at the right end of its interval than at the left end.

\paragraph{Step 10: the paper's example.} In the example, the incumbent is uniform on \([0,r]\), and \(J\) rises on
the interval \([r_0,r_2]\) of strengths (Chapter~\ref{ch:coexistence}). Proposition~3(c) gives that
\(\DT\,(g_H-c_H)\) is nondecreasing there, with \(g_H>c_H\), while \(\gamma=\bar\Gamma-\DT\) falls strictly where
\(\DT\) rises. A positive numerator that does not fall, over a positive denominator that falls, gives a ratio that rises.

\begin{checkyourself}\label{q:trading-negative}
On a path, suppose \(g_L(F_r)>c_H\) at the strength \(r\), and that \(g_L\) is the same at a stronger incumbent
\(r'\) while the spread widens strictly. What is the sign of \(\tau(r)\)? What is entry after each order at \(r\)?
Use the increment formula to show that \(\tau(r')<\tau(r)\).
\end{checkyourself}

% =====================================================================
\section{The race between the spread and the threshold}\label{sec:trading-race}

Put Lemma A.5 next to Proposition~1(b). Along a path, while \(g_L\le c_H\) and the profit gap stays positive, a
stronger incumbent does two things.
\begin{itemize}
\item It never narrows the spread. A wider spread raises what an informed order earns at every posterior and level
of entry, and in particular what a purchase that recruits earns.
\item It never lowers the threshold. A higher threshold means that a purchase must push the posterior further
before a high-cost challenger prepares, so recruitment becomes harder.
\end{itemize}
The two numbers pull in opposite directions. When \(\DT>0\) and \(g_H>g_L\), Lemma A.4 says that they are all that
the market sees of the auction.
So every question about orders, posteriors and entry along a path becomes a question about this race.

The race has no fixed winner: the return \(J\) of a purchase that just recruits falls in the instance of
Lemma A.5 and rises in the paper's example. The later results say which side wins under their own conditions.
\begin{itemize}
\item Chapter~\ref{ch:ccc} compares three incumbents. From a weak incumbent to a stronger one, the spread wins:
informed trade starts, and a purchase recruits a high-cost challenger. Against a still stronger incumbent, the
threshold wins: informed trade stays, but no order recruits.
\item Chapter~\ref{ch:coexistence} follows the race through the region where quiet and informative markets coexist.
The economy can enter that region, leave it and enter it again as the incumbent grows stronger.
\item Chapter~\ref{ch:research} adds research. A higher threshold needs more research before a purchase recruits,
and a wider spread raises what that purchase earns.
\end{itemize}
Each of these comparisons solves the equilibrium again in each economy. The deviation payoffs of this chapter, at
fixed prices and entry, are the tool inside each economy.

\begin{chaptersummary}
\begin{itemize}
\item Market makers put the expected entry gain into the price, so the investor earns only on the class: the gross
return is entry times the spread times the residual uncertainty, \(K_H=e(1)\DT[1-\mu(1)]\) and
\(K_L=e(-1)\DT\,\mu(-1)\).
\item The proceeds \(t_0\) and the low-class entry gain cancel, because the price is the market's average of the
same payoff that the investor receives.
\item A wrong-signed order earns at most \(-k\). Each class trades when its gross return exceeds \(k\), abstains when
it falls short, and may mix at equality under the default tie rules.
\item The trading bound \(\rho\,\kappa(\pi)\DT\le K_H,K_L\le\DT\) holds in every candidate, mixed or pure. Its factor
\(\kappa(\pi)=m(\min\{\pi,1-\pi\})\) peaks at \(1/(2+u)\) at the prior \(\tfrac12\), and the favored class binds.
\item The problem splits by class, because each return uses only the posterior and entry after its own order.
\item With \(\DT>0\) and \(g_H>g_L\), orders, posteriors, entry and research depend on the auction only through
\(\DT\) and \(\tau\); prices add \(t_0\) and \(t_L\), and ownership adds who wins.
\item A stronger incumbent never narrows the spread and, while \(g_L\le c_H\), never lowers the threshold. The two
numbers race, and the later chapters study who wins.
\end{itemize}
\end{chaptersummary}

% =====================================================================
\answersection

\answer{\ref{q:trading-level}}{Write the raised levels with primes: \(t_0'=t_0+z\), \(t_H'=t_H+z\),
\(t_L'=t_L+z\). The spread is unchanged, \(\DT'=t_H'-t_L'=\DT\), and so is the low-class gain,
\(t_L'-t_0'=t_L-t_0\). At the same posterior and entry,
\begin{align*}
P'(x) &= t_0'+e(x)\,\bigl[t_L'-t_0'+\DT'\mu(x)\bigr] \why{price formula}\\
&= t_0+z+e(x)\,\bigl[t_L-t_0+\DT\mu(x)\bigr] \why{\(t_0'=t_0+z\); gain and spread unchanged}\\
&= P(x)+z \why{price formula before the change}
\end{align*}
So every price rises by \(z\). The gross returns \(K_H=e(1)\DT[1-\mu(1)]\) and \(K_L=e(-1)\DT\mu(-1)\) contain only
entry, the spread and the posteriors, so they do not change. In words: the share pays \(z\) more in every outcome,
the price rises by \(z\), and the investor receives the extra \(z\) and pays it, as with \(t_0\). Section 2.2 of the
paper makes the same point for a common standalone component: it shifts prices and payoffs by a constant and leaves
incentives unchanged.}

\answer{\ref{q:trading-reveal}}{With \(\mu(1)=1\), \(K_H=e(1)\,\DT\,(1-1)=0\), so \(U_H(1)=-k<0=U_H(0)\), and class
\(H\) strictly prefers abstention. Class \(L\) does not buy either, because a purchase is wrong-signed for it and
earns at most \(-k\) (part (a)). So no informed investor buys. Then a purchase comes only from liquidity traders,
with probability \(b\) in each class, and Bayes' rule gives \(\mu(1)=\pi<1\), which contradicts \(\mu(1)=1\). The
slot also rules the situation out directly. A liquidity purchase arrives with probability \(b>0\) in both classes,
so the likelihood ratio of a purchase is at most \(1+u\). Hence \(\mu(1)\le M(\pi)<1\) under every profile
(Lemma A.1(b)).}

\answer{\ref{q:trading-prior}}{Since \(\pi>\tfrac12\), \(1-\pi<\tfrac12<\pi\). Then
\begin{align*}
1-M(\pi) &= m(1-\pi) \why{fact 1}\\
&< m(\pi) \why{fact 2 and \(1-\pi<\pi\)}
\end{align*}
So \(\kappa(\pi)=1-M(\pi)\), the floor of class \(H\): the investor in class \(H\) has the smaller guaranteed
return. The prior already leans toward class \(H\), and a purchase can push the posterior up to \(M(\pi)\), close to
one. Little weight on class \(L\) then remains in the price, and class \(H\) earns only on that weight. A sale by
class \(L\) leaves a posterior of at least \(m(\pi)\), which is larger, so the residual weight on class \(H\) that
class \(L\) earns on is larger.}

\answer{\ref{q:trading-levels-two}}{By Lemma A.4, the orders \((\sigma_H,\sigma_L)\), the posteriors \(\mu\) and
entry \(e\) agree, because they solve the same common conditions; with research, \(\hat\lambda\) agrees too. The
belief at the price of each order is \(\mu(x)\) in both auctions, and the high-cost preparation probability at that
price is \((e(x)-\rho)/(1-\rho)\) in both. The prices can differ. With primes for the second auction, and the same
\(e\), \(\mu\) and \(\DT\),
\begin{align*}
P(x)-P'(x) &= t_0+e(x)\,[t_L-t_0+\DT\mu(x)]-t_0'-e(x)\,[t_L'-t_0'+\DT\mu(x)] \why{price formulas}\\
&= (t_0-t_0')+e(x)\,\bigl[(t_L-t_0)-(t_L'-t_0')\bigr] \why{\(e(x)\DT\mu(x)\) cancels}
\end{align*}
Ownership can differ too, because it carries class \(H\)'s chance to win against each auction's incumbent, \(W_H(F)\).}

\answer{\ref{q:trading-negative}}{Here \(\tau(r)=[c_H-g_L(F_r)]/\gamma(r)\) has a negative numerator and a positive
denominator, so \(\tau(r)<0\). Gross profit \(B_r(\mu)=g_L(F_r)+\mu\gamma(r)\) is at least \(g_L(F_r)>c_H\) at every
belief \(\mu\ge0\), so a high-cost challenger prepares after every price, and entry is one after every order. The
question assumes that the threshold is defined at \(r'\), that is, \(\gamma(r')>0\). Then
\begin{align*}
\tau(r')-\tau(r) &= \frac{g_L(F_r)-g_L(F_{r'})+\tau(r)\,[\DT(r')-\DT(r)]}{\gamma(r')} \why{Lemma A.5}\\
&= \frac{\tau(r)\,[\DT(r')-\DT(r)]}{\gamma(r')}<0 \why{\(g_L\) unchanged; \(\tau(r)<0\), \(\DT(r')>\DT(r)\)}
\end{align*}
So the threshold falls. Entry does not change, because every posterior already lies above a negative threshold.}
